\section{Introduction}\label{sec:introduction}

We determine the leading small positive two-form eigenvalues of the
Ricci-flat metrics in a polarized Calabi--Yau conifold transition.
Their number is the loss of second cohomology under smoothing. Their
coefficients depend on the periods of divisor classes over exceptional
curves and on the global polarized intersection form. The same period
map determines the integral relations among the vanishing cycles.
Thus the topological and spectral conclusions use the same marked
geometric transition. Local node types alone do not determine either
the compact cohomology ring or these spectral coefficients.

We calculate the global data for the split quintic and for two toric
transitions constructed in earlier work. The two toric applications
use established geometry from the stability--degeneration series and
from the separate Calabi--Yau smoothing series. We denote Papers
IV--VI of the latter series by CY~IV--VI below; these are distinct
from the numbered stability--degeneration papers. The spectral theorem
and its proof come first, followed by the compact applications and the
precise roles of their geometric sources.

\subsection{Geometric hypotheses and the main theorem}
Let $X_0$ be a normal projective threefold with trivial dualizing sheaf,
with exactly $k\geq1$ ordinary double points and no other singularities.
Let
\[
 q:Y\longrightarrow X_0,\qquad
 f:\mathcal X\longrightarrow\Delta
\]
be respectively a projective small crepant resolution and a proper flat
polarized smoothing over a complex disc. Assume that $Y$ is smooth and
simply connected, that the relative canonical bundle of $f$ is trivial,
and that $X_t=f^{-1}(t)$ is smooth for $0<|t|<t_0$. At each node the local
smoothing function $\tau_i$ is required to satisfy
\begin{equation}\label{eq:simple-smoothing}
 \tau_i(0)=0,\qquad \tau_i'(0)\ne0.
\end{equation}
This is a condition on a genuine family over a disc. Let $C_i\simeq\PP^1$
be the exceptional curves and impose
\begin{equation}\label{eq:betti-hypotheses}
 b_2(Y)=2,\qquad b_2(X_t)=1.
\end{equation}
The polarization on $X_0$, its pullback to $Y$, and its extension to $X_t$
are denoted by $H$. It is an integral class, with no implicit factor of
$2\pi$. Choose an integral divisor class $D$ independent of $H$ over
$\RR$, and put
\begin{equation}\label{eq:intersection-data}
 v=H^3,\qquad a=DH^2,\qquad d=D^2H,\qquad
 b_i=D\cdot C_i,\qquad M=\frac{a^2}{v}-d.
\end{equation}
All intersections in this display are taken on $Y$.

On $X_t$ let $g_t$ be the unique Ricci-flat K\"ahler metric with
$[\omega_t]=H$. Let $g_0$ be the singular Ricci-flat K\"ahler metric
in class $H$ on $M_0=X_0^{\reg}$, supplied by \cite{egz,hein-sun}.
By central harmonic representatives we mean smooth $L^2$ two-forms
on $(M_0,g_0)$ that are closed and coclosed. Sections~\ref{sec:mass}
and~\ref{sec:multiple} construct the representatives used below and
justify their domains and uniqueness. We use the nonnegative real Hodge Laplacian
$\Delta_p=dd^*+d^*d$ on real $p$-forms and its ordinary $L^2$ norm.
Write $\lambda_2^+(t)$ for its least strictly positive eigenvalue in
degree two. The subscript indicates form degree. All limits below are
along the fixed positive ray $0<t<t_0$.

\begin{theorem}\label{thm:main}
Under the preceding hypotheses, $M>0$, and
\begin{equation}\label{eq:main-law}
 \lambda_2^+(t)=
 \frac{2\pi\sum_{i=1}^k b_i^2}{M\,|\log t|}(1+o(1)).
\end{equation}
For all sufficiently small $t$, this eigenvalue has real multiplicity one,
and its real eigenspace is generated by a coexact primitive $(1,1)$-form.
There is a constant $\gamma>0$ such that the degree-one spectrum and every
other positive degree-two eigenvalue are at least $\gamma$.
Moreover,
\begin{equation}\label{eq:main-upper}
 \lambda_2^+(t)\le
 \frac{2\pi\sum_i b_i^2}{M\,|\log t|}+C|\log t|^{-2}.
\end{equation}
The constants refer to the fixed family, polarization, ray and local
identifications. No uniformity as the leading coefficients in
\eqref{eq:simple-smoothing} degenerate is asserted.
\end{theorem}

Replacing $D$ by $cD+eH$, with $c\ne0$, multiplies both $M$ and
$\sum b_i^2$ by $c^2$. Thus the coefficient belongs to the one-dimensional
divisor quotient modulo the surviving polarization. It is expressed in
finite global intersection and restriction data; the list of local nodes
alone does not determine it. The theorem gives a matching leading
asymptotic and a sharper one-sided upper estimate. It does not assert a
two-sided second-order remainder.

\subsection{The complete fixed-family cluster}
The same geometric problem allows several divisor directions to be lost.
Retain the preceding family, polarization and simple-smoothing hypotheses,
but replace \eqref{eq:betti-hypotheses} by
$b_2(Y)=b$, $b_2(X_t)=s$ and $m=b-s>0$. Set
\[
 V=H^2(Y;\RR),\qquad r(D)=(D\cdot C_i)_i,\qquad
 K=\ker r,\qquad \mathcal Q=V/K.
\]
The attachment sequence identifies $K$ with the surviving cohomology
$H^2(X_t;\RR)$. Define the forms
\begin{equation}\label{eq:intro-matrix-data}
\begin{split}
 G(D,E)&=-\int_YDEH+
     \frac{3(\int_YDH^2)(\int_YEH^2)}{2\int_YH^3},\\
 \mathcal E([D],[E])&=r(D)\cdot r(E).
\end{split}
\end{equation}
Let $\mathcal M$ be the quotient form induced by $G$ on $\mathcal Q$,
equivalently the minimum squared norm over each coset.
The theorem includes the positivity needed for this definition.
Write $\lambda^+_{2,j}(t)$ for the ordered positive real degree-two
eigenvalues, with multiplicity.

\begin{theorem}\label{thm:matrix}
The form $G$ is the positive-definite inner product of central harmonic
representatives, and $\mathcal E,\mathcal M$ are positive definite on the $m$-dimensional
space $\mathcal Q$. If $0<\mu_1\le\cdots\le\mu_m$ are the
generalized eigenvalues of $\mathcal E$ relative to $\mathcal M$, then
\begin{equation}\label{eq:matrix-law}
 \lambda^+_{2,j}(t)=\frac{2\pi\mu_j}{|\log t|}(1+o(1)),
 \qquad 1\le j\le m.
\end{equation}
The degree-one spectrum and $\lambda^+_{2,m+1}(t)$ have a uniform
positive lower bound. For small $t$, the entire $m$-dimensional positive
small spectral space consists of coexact primitive real $(1,1)$-forms.
A multiplicity-one $\mu_j$ gives an eventually simple eigenvalue.
These assertions concern the fixed family and ray; no uniformity in
degenerating node-scale ratios or two-sided next-order error is asserted.
\end{theorem}

The quotient removes all surviving harmonic directions, rather than
only $H$. It is calculated by the intersection Schur complement
\eqref{eq:mass-schur}. When $b=2$ and $s=1$, this gives the coefficient
of Theorem~\ref{thm:main}; its one-sided sharper upper bound is proved
separately. Section~\ref{sec:multiple} supplies the full kernel
projection and the ordered lower bounds needed for
Theorem~\ref{thm:matrix}.

\subsection{Analytic inputs and proof}
Singular Ricci-flat metric existence, conical regularity and the local
expansions for the actual smoothing metrics are supplied by
\cite{egz,hein-sun,benabida}. They do not determine the compact
small-eigenvalue coefficients. We combine those metric results with
conic spectral compactness and a sharp capacity calculation to prove
\eqref{eq:main-law} and the ordered matrix law \eqref{eq:matrix-law}.

Ann\'e--Takahashi describe the spectral limits in exact-cone surgery
in terms of the central operator, complete-cap harmonic forms and critical
extended solutions \cite{anne-takahashi}. We use their norm decomposition
for finitely many caps, with the critical terms excluded in physical
degrees one and two. The proof then has three steps. The restriction map
fixes the harmonic link periods, and a scalar correction computes the
central harmonic Gram form from intersections. On exactly conical models,
logarithmic cutoffs give upper bounds and cap capacity gives matching lower
bounds on every limiting spectral subspace. Dodziuk's comparison of Hodge
eigenvalues under equivalent metrics transfers the entire ordered cluster
to the actual Ricci-flat family \cite{dodziuk}.

The leading law requires only the $C^0$ condition stated in
Definition~\ref{def:C0-comparison}; its geometric verification is separate
from the model argument. The differentiated estimate of
Proposition~\ref{prop:metric-comparison} is used for the sharper rank-one
upper remainder. Projection off all surviving harmonic classes is performed
on the models, where the needed compact convergence follows from the fixed
central elliptic problem.

\subsection{Three compact applications}
For independent linear forms $l_1,l_2$ and general quartics $q_1,q_2$ on
$\PP^4$, the split quintic $X_0=\{l_1q_2-l_2q_1=0\}$ has sixteen
nodes on the plane $\{l_1=l_2=0\}$. Its incidence resolution lies in
$\PP^4\times\PP^1$, with hyperplane classes $H,L$.

\begin{corollary}\label{cor:quintic}
Let $X_t=\{l_1q_2-l_2q_1+tG=0\}$, where the quintic $G$ is nonzero
at every node. For the Ricci-flat metrics in class $H$,
\[
 \lambda_2^+(t)\sim\frac{10\pi}{|\log t|}.
\]
The eigenvalue has all the simplicity, type, gap and one-sided remainder
properties of Theorem~\ref{thm:main}.
\end{corollary}

Here $r(L)=(1,\ldots,1)$ and $M=16/5$. The same primitive diagonal
vector generates the integral relation among the cohomological vanishing
cycles. The marked transition, including its full integral rings and
monodromy image of index sixteen, is stated in
Theorem~\ref{thm:integrated} and proved in Appendices~\ref{sec:integral-topology}
and~\ref{sec:tangent}.

\subsubsection*{Eight nodes}
Let $B=\PP^1\times\PP^1$ and
$Z=\PP_B(\OO\oplus\OO(1,-1))\times_B
\PP_B(\OO\oplus\OO(1,-1))$. In the quotient convention write
$f,g$ for the base hyperplane classes and $\xi_1,\xi_2$ for the
tautological classes. The anticanonical contraction of a general
anticanonical hypersurface $Y\subset Z$ is the eight-node threefold
$X_0$ constructed in \cite[Appendix~A, Proposition~A.1]{wuebben-chambers}.

\begin{corollary}\label{cor:eight-node}
This transition admits polarized smoothings with trivial relative
canonical bundle and first-order smoothing at every node, whose central
polarization pulls back to
$H=(4g+2\xi_1+2\xi_2)|_Y$. For every such fixed family and ray,
the actual Ricci-flat metrics in class $H$ have precisely two small
positive degree-two eigenvalues, with
\begin{equation}\label{eq:eight-spectral-law}
 \lambda^+_{2,1}(t)\sim\frac{\pi}{|\log t|},\qquad
 \lambda^+_{2,2}(t)\sim\frac{5\pi}{4|\log t|}.
\end{equation}
Both are eventually simple, coexact and primitive real $(1,1)$.
The degree-one spectrum and all other positive degree-two eigenvalues
have a uniform positive lower bound for that fixed family.
\end{corollary}

Here $b_2(Y)=4$ and $b_2(X_t)=2$. The period matrix and the quotient
mass are calculated from the earlier proposition's divisor intersections
and two exceptional-curve types. Section~\ref{sec:eight-node} verifies
the additional polarized-family hypotheses and proves the corollary.
This is an application of Theorem~\ref{thm:matrix}; it makes no new
integral saturation assertion for the eight-node family.

\subsubsection*{Four nodes}
Let $X'_9$ be the four-node partial resolution on the nonmirror side of
\cite[Section~11.1]{wuebben-cy-mirror}, with the chosen projective
small resolution $Y$. Section~\ref{sec:cy-x9} specifies the ample
Cartier class $\mathcal L_0$ by the eleven-ray support vector
\eqref{eq:x9-height}. Write $H=c_1(q^*\mathcal L_0)$.

\begin{corollary}\label{cor:x9}
The threefold $X'_9$ admits a relatively polarized smoothing carrying
exactly $\mathcal L_0$, with trivial relative canonical bundle and
first-order smoothing at every node. One has $b_2(Y)=6$,
$b_2(X_t)=3$ and $H^3=8712$. For every such fixed family and ray,
the actual Ricci-flat metrics in class $H$ have precisely three small
positive degree-two eigenvalues, with
\begin{equation}\label{eq:x9-spectral-law}
\begin{split}
 |\log t|\lambda^+_{2,1}(t)&\longrightarrow\frac{2\pi}{7},\\
 |\log t|\lambda^+_{2,2}(t)&\longrightarrow
       \frac{\pi(375-3\sqrt{3065})}{628},\\
 |\log t|\lambda^+_{2,3}(t)&\longrightarrow
       \frac{\pi(375+3\sqrt{3065})}{628}.
\end{split}
\end{equation}
All three are eventually simple, coexact and primitive real $(1,1)$.
The degree-one spectrum and all other positive degree-two eigenvalues
have a uniform positive lower bound for that fixed family.
\end{corollary}

The source's integral exceptional relation $(1,-1,-1,1)$ supplies a
full-support deformation tangent through
\cite[Theorems~2.5 and~4.5]{wuebben-cy-smoothing}. We prove that the
exact ample class extends, and calculate the full three-dimensional
harmonic quotient. The resulting spectral law is an additional
consequence on the established original-side geometry; it is not a
claim about the mixed-singularity $X_9$ or the distinct mirror family.

\subsection{Relation to earlier work}
The split-quintic construction and conifold attachment sequences are
classical \cite{greene-morrison-strominger,rossi,lee-lin-wang}.
Our two toric applications use specific earlier geometric constructions:
\cite[Appendix~A, Proposition~A.1]{wuebben-chambers} supplies the eight-node
contraction and Picard rank, while \cite[Section~11.1]{wuebben-cy-mirror}
supplies the four-node threefold $X'_9$ and its exceptional circuits.
The nodal specialization of \cite[Theorems~2.5 and~4.5]{wuebben-cy-smoothing}
provides the latter's deformation tangent and smooth deformation base.
In both applications we check first-order smoothing and extension of the
specified ample class before applying the spectral theorem.

The curve-class map and the divisor-period map are transposes:
\begin{equation}\label{eq:curve-period-duality}
 A:\QQ^k\longrightarrow H_2(Y;\QQ),\quad e_i\longmapsto[C_i],
 \qquad r=A^*,\qquad \im r=(\ker A)^\perp.
\end{equation}
The space $\ker A$ governs localized deformation tangents, whereas $\ker r$
is the surviving second cohomology and $\im r$ supplies the spectral
periods. The harmonic quotient also requires the polarized intersection
form. Integral saturation is proved from the actual attachments where it
is asserted; it does not follow from \eqref{eq:curve-period-duality}.

Ann\'e--Takahashi's conic surgery theorem includes logarithmic modes
\cite{anne-takahashi}. Rowlett proves spectral and heat convergence under
her asymptotically conic convergence hypotheses \cite{rowlett}.
Guillarmou--Sher's low-energy resolvent theorem \cite[Theorem~1]{guillarmou-sher}
excludes the harmonic link two-form channel present here through its
condition~(4). For scalar Laplacians in different degenerations, Cao proves
two-sided inverse-logarithmic orders for induced total-space metrics
\cite[Theorems~0.2--0.4]{cao}, and Dai--Yoshikawa obtain small-eigenvalue
orders and the leading product in reduced curve degenerations
\cite[Theorems~0.2--0.3]{dai-yoshikawa}. These results concern different
operators and metrics. Here the conclusion is the individual leading
coefficient of every lost-divisor mode for the actual polarized Ricci-flat
metrics, computed from the full quotient harmonic norm.

The tangent connections in Appendix~\ref{sec:tangent} are induced by those
Ricci-flat metrics. Collins--Picard--Yau treat tangent-bundle HYM geometry
on balanced conifold transitions \cite{collins-picard-yau}; our appendix
identifies the flat tangent-sheaf limit and its reflexive hull for the
specified K\"ahler family.

Section~\ref{sec:geometry} establishes the actual metric comparison.
Section~\ref{sec:mass} calculates the harmonic norm and the integral
restriction map. Sections~\ref{sec:capacity}--\ref{sec:multiple} prove the
model spectral law, transfer it using only $C^0$ comparison, and isolate
the stronger rank-one remainder. Sections~\ref{sec:quintic},
\ref{sec:eight-node} and~\ref{sec:cy-x9} give the three applications.
The appendices retain the supporting integral and tangent calculations.
